\chapter{空间映射估计}
\label{chap:spatial-map-estimation}

测量设备常常只给出相对的几何信息。测距仪报告两个标记点相隔多远，却不报告它们在哪个坐标系中；两台相机拍到同一组标记，各自记下自己的像素坐标，却不知道两套坐标之间怎样换算。这里有两类问题：只知道点与点之间的距离，能否给这些点安排坐标；知道两组点一一对应，能否求出把一组坐标送到另一组坐标的映射。两类问题都要先问观测究竟确定了什么，再谈怎样计算。

本章沿“可辨识性—坐标重构—映射估计”的顺序展开。第~\ref{sec:smap-identifiability}节说明距离对整体平移、旋转和反射一无所知，缺失距离还可能留下真正的形变自由度，固定参照只能消除前一类自由度。第~\ref{sec:smap-distance-reconstruction}节处理距离表：完整表经过平方和双中心化变成Gram矩阵，由谱分解读出能否实现、最少需要几维以及坐标本身；再研究指定维数下的近似、只观测部分距离时形状何时确定，以及测量误差怎样传到坐标。

第~\ref{sec:smap-map-estimation}节转向已知对应点的映射估计。先规定允许的映射族，再依次处理刚体、仿射和射影三种模型：刚体模型由质心消去平移、由交叉矩阵的奇异值分解求旋转；仿射模型用增广坐标化为线性最小二乘；射影模型在齐次坐标下只确定到一个非零倍数，代数残差与几何上的重投影误差需要分开评价。每种模型都逐项核对数据是否可满足、参数是否唯一、扰动下是否稳定。

本章把点之间的对应关系当作给定的输入，不讨论怎样从未标注的两组点中找出对应。旋转、位姿与射影空间的定义取自第~\ref{chap:lie-groups}章；以几何误差为目标的迭代求解属于卷五的光滑优化。有限个点的另一类问题，即把样本看成某个未知流形上的采样、恢复它的局部几何与拓扑，是第~\ref{chap:manifold-recovery}章的任务。

\section{相对观测的可辨识性}
\label{sec:smap-identifiability}

在计算任何坐标或映射之前，先要弄清观测能够确定什么。一组测量即使完全没有误差，也可能对某些量毫无约束：把整幅图平移一下，所有距离都不变，距离表当然无法告诉我们原点在哪里。本节先区分两类相对观测，再把观测不能辨认的自由度分成两种：一种来自坐标系的任意选择，另一种来自观测本身的缺失。

\paragraph{两类相对观测}

第一类观测只涉及点与点之间的关系。给定点标签集合$V=\{1,\ldots,n\}$和已观测的点对集合$E$，对每个$\{i,j\}\in E$测得一个距离$d_{i,j}\geq0$。未知的是坐标$\bm x_1,\ldots,\bm x_n\in\mathbb R^d$，要求它们满足
\[
 \lVert\bm x_i-\bm x_j\rVert_2=d_{i,j},\qquad\{i,j\}\in E .
\]
$E$包含全部点对时称距离表完整，否则称为部分距离。这里的数据只有标签与数，没有任何外部坐标。

第二类观测涉及同一组物理点在两种坐标表示下的记录。对每个$i$，已知源坐标$\bm x_i\in\mathbb R^d$与目标坐标$\bm y_i\in\mathbb R^m$，并且知道它们描述同一个点。未知的是把源坐标送到目标坐标的映射$f$，要求$f(\bm x_i)\approx\bm y_i$。两类问题的未知量不同：前者的未知量是$n$个点的坐标，共$nd$个数；后者的未知量是一个映射，点越多，对它的约束越多。

\paragraph{整体欧氏自由度}

先看距离不能辨认什么。取正交矩阵$\bm Q\in\mathrm O(d)$和向量$\bm b\in\mathbb R^d$，令$\bm x_i'=\bm Q\bm x_i+\bm b$。差向量变为$\bm x_i'-\bm x_j'=\bm Q(\bm x_i-\bm x_j)$，正交矩阵保持长度，所以每一对点的距离都不变。也就是说，平移、旋转和反射作用于整组点以后，任何距离观测都看不出区别。

\begin{definition}[合同的点配置]
\label{def:smap-congruent}
两组点$\bm x_1,\ldots,\bm x_n$与$\bm x_1',\ldots,\bm x_n'\in\mathbb R^d$称为\term{合同}（congruent），如果存在$\bm Q\in\mathrm O(d)$和$\bm b\in\mathbb R^d$，使$\bm x_i'=\bm Q\bm x_i+\bm b$对所有$i$成立。
\end{definition}

由上一段的计算，合同的配置给出同一张距离表。因此距离观测最多只能把点配置确定到合同的程度，这一限制与测量是否完整、是否精确都无关。平移有$d$个参数，$\mathrm O(d)$的每个连通分支有$d(d-1)/2$个参数，所以合同带来$d(d+1)/2$个连续自由度，外加一个反射的二选一。平面上是$3$个连续自由度，空间中是$6$个。

反过来，完整距离表是否把配置确定到合同为止？第~\ref{sec:smap-full-distance}节会证明答案是肯定的：两组点的全部距离相同，它们就合同。这说明对完整距离表而言，合同是唯一的不确定性，坐标重构的任务是在合同意义下找出配置。

\paragraph{缺失观测的形变自由度}

观测不完整时，不确定性可能多于合同。在平面上取单位正方形的四个顶点，只测四条边长，都等于$1$，不测两条对角线。把正方形压成内角为$\theta$的菱形，四条边长仍是$1$，但两条对角线长分别变为$2\sin(\theta/2)$与$2\cos(\theta/2)$。$\theta\neq\pi/2$时对角线与正方形不同，所以菱形与正方形不合同，却满足全部观测。

这里的自由度来自观测不足，可以用计数来理解。平面上$n$个点有$2n$个坐标，扣去合同的$3$个自由度，形状本身还有$2n-3$个自由度；$n=4$时为$5$个，四个边长约束一般只能消去其中四个，剩下的一个就是菱形的内角。计数只给出必要条件，约束够多也未必确定形状，第~\ref{sec:smap-partial-distance}节将用具体例子说明这一点。

\paragraph{固定参照的作用}

为了得到一组具体坐标，常常人为固定参照。在平面上可以规定$\bm x_1=\bm0$，$\bm x_2$位于第一坐标轴正半轴上，$\bm x_3$位于上半平面。这三条规定恰好消去平移、旋转和反射：任何配置都与唯一一个满足规定的配置合同，只要$\bm x_1,\bm x_2,\bm x_3$不共线。

但规定只能消去约定上的自由度。回到菱形的例子，取$\bm x_1=\bm0$、$\bm x_2=(1,0)^\top$，并把第四个点放在上半平面，内角$\theta$仍然可以连续变化，四个边长仍然成立。形变自由度是由缺失的观测造成的，再多的坐标约定也补不出缺失的对角线长度。

还要区分约定与观测。若某些点的绝对位置是用仪器测出来的，例如卫星定位给出的锚点坐标，它们就是额外的观测，可能与距离数据相互矛盾，需要进入拟合残差；只有任意选取、不影响距离的规定，才是消去合同自由度的约定。第二类问题中也有对应的区分：两套坐标系都已给定，映射本身没有约定上的自由度，但必须事先选定映射族，第~\ref{sec:smap-correspondence-models}节将说明这一点。

\section{距离观测的坐标重构}
\label{sec:smap-distance-reconstruction}

上一节说明，距离最多把点配置确定到合同为止。本节在这个限度内研究重构：完整距离表在什么条件下来自欧氏点、最少需要几维、怎样写出坐标；维数受限时怎样近似；只观测部分距离时形状是否仍然确定；以及测量误差会怎样传到坐标。比较两组坐标时，总要先消去整体平移与正交变换，否则相差一个旋转的两组正确坐标也会被误判为不同。

\subsection{完整距离的欧氏实现}
\label{sec:smap-full-distance}

设$n$个点两两之间的距离全部已知，排成矩阵$\bm D=[d_{i,j}]\in\mathbb R^{n\times n}$。任何来自点的距离表都满足三条基本要求：各项非负，$d_{i,j}=d_{j,i}$，对角元$d_{i,i}=0$。下面始终假定输入满足这三条，并记逐项平方后的矩阵为$\bm\Delta=[\delta_{i,j}]$，$\delta_{i,j}=d_{i,j}^2$。问题是：是否存在$\bm x_1,\ldots,\bm x_n\in\mathbb R^d$，使$\lVert\bm x_i-\bm x_j\rVert_2^2=\delta_{i,j}$对所有$i,j$成立？若存在，称$\bm\Delta$\term{欧氏可实现}（Euclidean realizable），所取的点称为它的一个实现。

\paragraph{从坐标到Gram矩阵}

先从已知坐标出发，看平方距离与内积怎样联系。把各点按行排成$\bm X\in\mathbb R^{n\times d}$，第$i$行为$\bm x_i^\top$，Gram矩阵$\bm G=\bm X\bm X^\top$记录两两内积$g_{i,j}=\bm x_i^\top\bm x_j$。展开平方距离，
\begin{equation}
 \delta_{i,j}=\lVert\bm x_i\rVert_2^2+\lVert\bm x_j\rVert_2^2-2\bm x_i^\top\bm x_j=g_{i,i}+g_{j,j}-2g_{i,j}.
 \label{eq:smap-distance-gram}
\end{equation}
写成矩阵，记$\bm g=(g_{1,1},\ldots,g_{n,n})^\top$、$\bm1$为全一向量，就是$\bm\Delta=\bm g\bm1^\top+\bm1\bm g^\top-2\bm G$。

反方向不能直接解出$\bm G$。距离不依赖原点，内积却依赖：整组点平移后，$\bm\Delta$不变而$\bm G$改变。所以要先约定原点，最自然的选择是质心。令
\begin{equation}
 \bm J=\bm I_n-\frac1n\bm1\bm1^\top ,
 \label{eq:smap-centering-matrix}
\end{equation}
它对称，满足$\bm J^2=\bm J$与$\bm J\bm1=\bm0$，是到$\bm1$的正交补上的正交投影。$\bm J\bm X$从每一行减去各点的平均坐标，所以$\bm J\bm X$的各行就是以质心为原点的坐标。

用$\bm J$左右夹乘$\bm\Delta$。由$\bm J\bm1=\bm0$，$\bm J\bm g\bm1^\top\bm J$与$\bm J\bm1\bm g^\top\bm J$都是零矩阵，只剩
\begin{equation}
 -\frac12\bm J\bm\Delta\bm J=\bm J\bm G\bm J=(\bm J\bm X)(\bm J\bm X)^\top .
 \label{eq:smap-double-centering}
\end{equation}
右端正是中心化坐标的Gram矩阵。这一运算称为\term{双中心化}（double centering）：左乘$\bm J$消去行上的均值项，右乘消去列上的均值项，把平方距离转成以质心为原点的内积。下面记$\bm B=-\frac12\bm J\bm\Delta\bm J$，它对任何对称零对角的$\bm\Delta$都有定义，不必事先知道$\bm\Delta$来自点。

\paragraph{可实现的判据}

若$\bm\Delta$可实现，由式~\eqref{eq:smap-double-centering}，$\bm B$是一个Gram矩阵，因而半正定：$\bm c^\top\bm B\bm c=\lVert(\bm J\bm X)^\top\bm c\rVert_2^2\geq0$。要证明半正定也是充分的，需要确认双中心化没有丢掉信息，即零对角的对称矩阵由它的双中心化唯一确定。

\begin{lemma}[双中心化的核]
\label{lem:smap-double-centering-kernel}
设$\bm H\in\mathbb R^{n\times n}$对称且$\bm J\bm H\bm J=\bm0$，则存在$\bm a\in\mathbb R^n$使$\bm H=\bm a\bm1^\top+\bm1\bm a^\top$。若还有$h_{i,i}=0$对所有$i$成立，则$\bm H=\bm0$。
\end{lemma}
\begin{proof}
记$\bm m=\frac1n\bm H\bm1$，$\mu=\frac1{n^2}\bm1^\top\bm H\bm1$。把$\bm J=\bm I-\frac1n\bm1\bm1^\top$代入$\bm J\bm H\bm J=\bm0$并展开，
\[
 \bm H-\bm1\bm m^\top-\bm m\bm1^\top+\mu\bm1\bm1^\top=\bm0 .
\]
取$\bm a=\bm m-\frac\mu2\bm1$，即得$\bm H=\bm a\bm1^\top+\bm1\bm a^\top$。此时$h_{i,i}=2a_i$，零对角迫使$\bm a=\bm0$，于是$\bm H=\bm0$。
\end{proof}

\begin{theorem}[欧氏可实现的判据]
\label{thm:smap-realizability}
设$\bm\Delta\in\mathbb R^{n\times n}$对称、零对角，$\bm B=-\frac12\bm J\bm\Delta\bm J$。则$\bm\Delta$是某组欧氏点的平方距离矩阵，当且仅当$\bm B$半正定。此时实现所需的最小维数等于$r=\operatorname{rank}\bm B$，并且可以取
\begin{equation}
 \bm Y=\bm Q_r\bm\Lambda_r^{1/2}\in\mathbb R^{n\times r},
 \label{eq:smap-classical-scaling}
\end{equation}
其中$\bm\Lambda_r$是$\bm B$的$r$个正特征值组成的对角矩阵，$\bm Q_r$的列是对应的标准正交特征向量。
\end{theorem}
\begin{proof}
必要性已在上面说明。设$\bm B$半正定，由对称矩阵的谱分解，$\bm B=\bm Q_r\bm\Lambda_r\bm Q_r^\top$，于是式~\eqref{eq:smap-classical-scaling}的$\bm Y$满足$\bm Y\bm Y^\top=\bm B$。由$\bm B\bm1=\bm0$得$\lVert\bm Y^\top\bm1\rVert_2^2=\bm1^\top\bm B\bm1=0$，即$\bm Y$的各行已经中心化，$\bm J\bm Y=\bm Y$。

记$\bm\Delta'$为$\bm Y$各行的平方距离矩阵。由式~\eqref{eq:smap-double-centering}，$-\frac12\bm J\bm\Delta'\bm J=\bm Y\bm Y^\top=\bm B$，所以$\bm H=\bm\Delta-\bm\Delta'$对称、零对角且$\bm J\bm H\bm J=\bm0$。由引理~\ref{lem:smap-double-centering-kernel}，$\bm H=\bm0$，即$\bm Y$实现了$\bm\Delta$。

最后看维数。任何$\mathbb R^k$中的实现，中心化后的Gram矩阵就是$\bm B$，而它是$n\times k$矩阵与其转置之积，秩不超过$k$，故$k\geq r$；上面的构造又在$k=r$时达到。$r=0$时所有点重合。
\end{proof}

这一判据与谱分解构造合称\term{经典多维尺度}（classical multidimensional scaling）。所用的对称矩阵谱分解可参考\BookChapterRef{chap:eigen-spectral-decomposition}{卷二第4章“特征值与谱分解”}。计算步骤是：逐项平方得$\bm\Delta$，双中心化得$\bm B$，求$\bm B$的特征值与特征向量，检查有无负特征值，再按式~\eqref{eq:smap-classical-scaling}写出坐标，最后用重构坐标重新计算距离与原表核对。

\paragraph{坐标的唯一性}

谱分解中特征向量可以变号，重特征值时特征空间内部还可以换基，所以式~\eqref{eq:smap-classical-scaling}给出的坐标并不唯一。下面的定理说明，这些差别恰好是正交变换，并且任何两个同维实现都只差一个合同。

\begin{theorem}[完整距离确定合同类]
\label{thm:smap-full-distance-congruence}
设$\bm x_1,\ldots,\bm x_n$与$\bm x_1',\ldots,\bm x_n'\in\mathbb R^d$的两两距离全部相同，则两组点合同。更具体地，以各自质心为原点后，存在$\bm Q\in\mathrm O(d)$使$\bm x_i'-\overline{\bm x}'=\bm Q(\bm x_i-\overline{\bm x})$对所有$i$成立。
\end{theorem}
\begin{proof}
记$\bm z_i=\bm x_i-\overline{\bm x}$、$\bm z_i'=\bm x_i'-\overline{\bm x}'$。由式~\eqref{eq:smap-double-centering}，两组中心化坐标的Gram矩阵都等于同一个$\bm B$，即$\bm z_i^\top\bm z_j=\bm z_i'^\top\bm z_j'$对所有$i,j$成立。

于是对任意系数$\bm c\in\mathbb R^n$，$\lVert\sum_ic_i\bm z_i\rVert_2^2=\bm c^\top\bm B\bm c=\lVert\sum_ic_i\bm z_i'\rVert_2^2$。取$S=\operatorname{span}\{\bm z_i\}$、$S'=\operatorname{span}\{\bm z_i'\}$，规定$L(\sum_ic_i\bm z_i)=\sum_ic_i\bm z_i'$。若两组系数给出$S$中同一向量，它们的差$\bm c$使$\sum_ic_i\bm z_i=\bm0$，由上式$\sum_ic_i\bm z_i'=\bm0$，所以$L$是良定义的线性映射，并且保持长度。

保长的线性映射是单射，所以$\dim S'=\dim S$，正交补$S^\perp$与$S'^\perp$维数也相同。在两个正交补中各取一组标准正交基，把前者逐个送到后者，就把$L$延拓为$\mathbb R^d$上保长的线性映射$\bm Q$，即正交矩阵，并且$\bm Q\bm z_i=\bm z_i'$。
\end{proof}

这一定理兑现了第~\ref{sec:smap-identifiability}节的论断：对完整距离表，合同是唯一的不确定性。$\bm\Lambda_r$的特征值互不相同时，式~\eqref{eq:smap-classical-scaling}中特征向量的选取只差各列的符号，对应沿各坐标轴的反射；有重特征值时，相应的坐标轴可以在特征空间内任意旋转。

\paragraph{实现与不可实现的例子}

\begin{example}[三个点的重构]
\label{ex:smap-three-point-mds}
取$\bm x_1=(-1,0)^\top$、$\bm x_2=(0,1)^\top$、$\bm x_3=(1,0)^\top$，平方距离为$\delta_{1,2}=\delta_{2,3}=2$、$\delta_{1,3}=4$。记一般的$\delta_{1,2}=\delta_{2,3}=a$、$\delta_{1,3}=b$，由对称性可以直接读出$\bm B$的特征向量：$\bm1$对应特征值$0$；对$\bm v=(1,0,-1)^\top/\sqrt2$，由$\bm J\bm v=\bm v$得$\bm v^\top\bm B\bm v=-\frac12\bm v^\top\bm\Delta\bm v=b/2$；对$\bm w=(1,-2,1)^\top/\sqrt6$，同样算得$\bm w^\top\bm B\bm w=(4a-b)/6$。

代入$a=2$、$b=4$，正特征值为$2$与$2/3$，所以精确重构恰好需要二维。由式~\eqref{eq:smap-classical-scaling}，重构坐标的两列为$\sqrt2\,\bm v=(1,0,-1)^\top$与$\sqrt{2/3}\,\bm w=(1,-2,1)^\top/3$，三个点是$(1,1/3)$、$(0,-2/3)$、$(-1,1/3)$。它们两两的平方距离为$1+1=2$、$1+1=2$、$4+0=4$，与原表一致。原配置平移到质心后为$(-1,-1/3)$、$(0,2/3)$、$(1,-1/3)$，重构结果恰是它旋转$180^\circ$的结果，与原配置合同；特征向量取相反的符号，就得到中心化的原配置本身。
\end{example}

\begin{example}[违反三角不等式的距离表]
\label{ex:smap-nonrealizable}
取$d_{1,2}=d_{2,3}=1$、$d_{1,3}=3$，即$a=1$、$b=9$。按上例的公式，$\bm B$的特征值为$9/2$、$(4-9)/6=-5/6$与$0$，出现负特征值，所以这张表在任何维数的欧氏空间中都不可实现。这与直接观察一致：$3>1+1$违反三角不等式。负特征值的特征向量$(1,-2,1)^\top$指出问题出在中间点相对两端的位置上。
\end{example}

\subsection{低维坐标重构}
\label{sec:smap-low-dimension}

实际中常常指定维数$k$，例如要在平面上画出点的分布，而$\bm B$的秩可能大于$k$；含误差的距离表还可能让$\bm B$出现负特征值，根本不可实现。这时只能求近似坐标$\bm Y\in\mathbb R^{n\times k}$。近似必须先说明比较什么，本小节先以Gram矩阵为比较对象，再看距离上的误差。

\paragraph{Gram矩阵的最优截断}

以Frobenius范数比较Gram矩阵，问题是在秩不超过$k$的半正定矩阵中找离$\bm B$最近的一个。$\bm B$半正定时，它的谱分解就是奇异值分解，截断到前$k$个特征值即得最优解，这是最佳低秩逼近定理的直接推论，可参考\BookChapterRef{chap:matrix-approximation}{卷二第15章“矩阵逼近”}。$\bm B$有负特征值时还要多一步：负谱部分对任何半正定近似都只会增加误差。

\begin{theorem}[半正定低秩近似]
\label{thm:smap-psd-truncation}
设$\bm B$对称，特征值$\lambda_1\geq\cdots\geq\lambda_n$，特征向量为标准正交的$\bm q_1,\ldots,\bm q_n$。令$p$为正特征值的个数，$k'=\min(k,p)$，$\bm B_k=\sum_{l\leq k'}\lambda_l\bm q_l\bm q_l^\top$。则对一切秩不超过$k$的半正定矩阵$\bm X$，
\[
 \lVert\bm B-\bm X\rVert_{\mathrm F}^2\geq\lVert\bm B-\bm B_k\rVert_{\mathrm F}^2=\sum_{l=k'+1}^{p}\lambda_l^2+\sum_{\lambda_l<0}\lambda_l^2 .
\]
\end{theorem}
\begin{proof}
把$\bm B$分成正部与负部：
\begin{equation*}
\bm B_+=\sum_{\lambda_l>0}\lambda_l\bm q_l\bm q_l^\top,\qquad\bm B_-=\sum_{\lambda_l<0}(-\lambda_l)\bm q_l\bm q_l^\top,
\end{equation*}
二者都半正定，$\bm B=\bm B_+-\bm B_-$，且由特征向量正交，$\bm B_+\bm B_-=\bm0$。对半正定的$\bm X$展开，
\[
 \lVert\bm B_+-\bm X-\bm B_-\rVert_{\mathrm F}^2
 =\lVert\bm B_+-\bm X\rVert_{\mathrm F}^2+\lVert\bm B_-\rVert_{\mathrm F}^2+2\operatorname{tr}(\bm X\bm B_-),
\]
这里用了$\operatorname{tr}(\bm B_+\bm B_-)=0$。两个半正定矩阵之积的迹非负：写$\bm B_-=\bm C\bm C^\top$，则$\operatorname{tr}(\bm X\bm B_-)=\operatorname{tr}(\bm C^\top\bm X\bm C)\geq0$。

所以$\lVert\bm B-\bm X\rVert_{\mathrm F}^2\geq\lVert\bm B_+-\bm X\rVert_{\mathrm F}^2+\lVert\bm B_-\rVert_{\mathrm F}^2$。$\bm B_+$半正定，由最佳低秩逼近定理，秩不超过$k$时$\lVert\bm B_+-\bm X\rVert_{\mathrm F}^2\geq\sum_{l=k'+1}^p\lambda_l^2$，等号在$\bm X=\bm B_k$时取到；而$\bm X=\bm B_k$时$\operatorname{tr}(\bm B_k\bm B_-)=0$，上面的不等式都取等号。
\end{proof}

于是低维坐标取$\bm Y_k=\bm Q_{k'}\bm\Lambda_{k'}^{1/2}$，即只保留前$k'$个正特征值对应的列。它仍然中心化，因为每个保留的特征向量都与$\bm1$正交。负特征值被舍弃，舍弃量$\sum_{\lambda_l<0}\lambda_l^2$应与测量误差的大小比较：若它远大于误差能解释的范围，问题出在数据或模型上，而不是维数选小了。

\paragraph{距离上的误差}

Gram矩阵的误差并不直接说明距离差了多少，需要单独计算。设$\bm\Delta$可实现，$\bm Y=\bm Q_r\bm\Lambda_r^{1/2}$是它的中心化实现。$\bm Y_k$恰好是$\bm Y$的前$k$列，相当于把每个点正交投影到前$k$个坐标轴张成的子空间上。正交投影不增加长度，所以每一对点的距离只会缩短：$\lVert\widetilde{\bm y}_i-\widetilde{\bm y}_j\rVert_2\leq d_{i,j}$。

缩短的总量可以精确算出。对中心化的点，$\sum_{i<j}\lVert\bm y_i-\bm y_j\rVert_2^2=n\sum_i\lVert\bm y_i\rVert_2^2=n\operatorname{tr}\bm B$，这只需展开平方并用$\sum_i\bm y_i=\bm0$。对截断坐标同样有$\sum_{i<j}\widetilde\delta_{i,j}=n\operatorname{tr}\bm B_k$，相减得
\[
 \sum_{i<j}\bigl(\delta_{i,j}-\widetilde\delta_{i,j}\bigr)=n\sum_{l>k}\lambda_l .
\]
注意这里是特征值之和，而Gram误差的平方是特征值的平方和。两种误差都由舍弃的谱决定，但不是同一个量。

\begin{example}[三个点的一维近似]
\label{ex:smap-one-dimensional}
回到例~\ref{ex:smap-three-point-mds}，取$k=1$，只保留特征值$2$。一维坐标为$(1,0,-1)$，两端距离仍为$2$，中间点到两端的距离由$\sqrt2$缩成$1$。Gram误差为舍弃的特征值$2/3$；平方距离的总缩短量为$(2-1)+(2-1)+0=2$，恰为$n\lambda_2=3\cdot2/3$。
\end{example}

\paragraph{最优性依赖目标}

谱截断在Gram矩阵的Frobenius误差下最优，不等于在一切几何准则下最优。常用的另一种准则直接比较距离，称为应力：$\sum_{i<j}(\lVert\widetilde{\bm y}_i-\widetilde{\bm y}_j\rVert_2-d_{i,j})^2$。在上例中，截断坐标的应力为$2(\sqrt2-1)^2=2(3-2\sqrt2)\approx0.343$。

改取一维坐标$(-s,0,s)$，应力为$2(s-\sqrt2)^2+(2s-2)^2$，对$s$求导令其为零，得$s=(2+\sqrt2)/3\approx1.138$。代回去，$s-\sqrt2=(2-2\sqrt2)/3$，$2s-2=(2\sqrt2-2)/3$，应力为$3\cdot4(\sqrt2-1)^2/9=4(3-2\sqrt2)/3\approx0.229$，比截断坐标小。截断坐标把两端拉得太远、中间缩得太多，而应力准则愿意把各对距离的误差摊平。应力关于坐标不是凸的，一般需要迭代优化，而不能由一次谱分解得到。

\subsection{部分距离的重构}
\label{sec:smap-partial-distance}

许多测量只能给出部分点对的距离，例如传感器只能测到通信范围内的邻居。把已观测的点对看作边，就得到\term{观测图}（observation graph）$\mathcal G=(V,E)$；图的基本概念可参考\BookChapterRef{chap:graph-objects}{卷三第1章“图与连接结构”}。观测图不连通时，不同分量之间没有任何约束，各分量可以独立地移动，这是显然的；下面两个例子说明，连通也远远不够。

\paragraph{连通而形状不定}

在平面上取三个点，只观测$d_{1,2}=d_{2,3}=1$。观测图是一条路径，是连通的。但$\bm x_2$处的夹角可以任意选取：$\bm x_1$与$\bm x_3$都在以$\bm x_2$为圆心的单位圆上，它们的距离可以是$[0,2]$中的任何值。不同的夹角给出互不合同的配置，所以连通不保证形状确定。

\paragraph{局部确定而全局不唯一}

为了区分两种“确定”，先给出定义。设$\bm x_1,\ldots,\bm x_n\in\mathbb R^d$满足观测图$\mathcal G$上的全部距离，称它是一个实现。

\begin{definition}[局部刚性与全局刚性]
\label{def:smap-rigidity}
若存在$\eta>0$，使得每个满足同一组观测距离、且$\lVert\bm x_i'-\bm x_i\rVert_2<\eta$对所有$i$成立的实现都与原实现合同，则称原实现\term{局部刚性}（locally rigid）。若满足观测距离的每个实现都与它合同，则称它\term{全局刚性}（globally rigid）。
\end{definition}

取平面上$\bm x_1=(0,0)^\top$、$\bm x_2=(2,0)^\top$、$\bm x_3=(1,1)^\top$、$\bm x_4=(1,-2)^\top$，观测除$\{3,4\}$外的五对距离：$d_{1,2}=2$，$d_{1,3}=d_{2,3}=\sqrt2$，$d_{1,4}=d_{2,4}=\sqrt5$。这是两个共用边$\{1,2\}$的三角形。

先看局部刚性。三角形的三边确定它到合同为止，所以$\bm x_1,\bm x_2,\bm x_3$整体只能作合同变换；在此基础上，$\bm x_4$到$\bm x_1,\bm x_2$的距离给定，它只能落在两个圆的两个交点之一，即直线$\bm x_1\bm x_2$的两侧各一个。两个交点相距$4$，小范围的扰动不能从一个跳到另一个，所以原实现局部刚性。

再看全局。把$\bm x_4$关于直线$\bm x_1\bm x_2$反射到$(1,2)^\top$，五个观测距离都不变，未观测的$d_{3,4}$却由$3$变成$1$。两个实现互不合同，所以原实现不是全局刚性的。局部确定只排除了附近的连续形变，不排除远处另有一个离散的解。

\paragraph{缺边、不相容与含噪}

固定维数$d$，并按第~\ref{sec:smap-identifiability}节的方式约定参照以消去合同自由度之后，部分距离问题有三种不同的困难。缺边可能使解不唯一，如上面的路径和两个三角形；观测不相容时根本无解，例如三个点的三条观测边长为$1$、$1$、$3$；观测含噪时一般也无精确解，只能在已观测的边上度量偏差。

含噪时自然的做法是只对已观测的边建立残差，例如
\[
 \sum_{\{i,j\}\in E}\bigl(\lVert\bm x_i-\bm x_j\rVert_2-d_{i,j}\bigr)^2 ,
\]
或用平方距离之差代替距离之差。未观测的点对不进入目标，因为对它们没有任何数据。这类目标关于坐标不是凸的，可能有多个局部极小，对应的正是上面的离散多解；其求解属于迭代优化。

\paragraph{不能把缺失填成零}

把缺失的距离填零，再套用第~\ref{sec:smap-full-distance}节的谱重构，是常见的错误。零距离不是“未知”，而是断言两点重合。在路径的例子中填$d_{1,3}=0$，得到完整表$\delta_{1,2}=\delta_{2,3}=1$、$\delta_{1,3}=0$，按例~\ref{ex:smap-three-point-mds}的公式，$a=1$、$b=0$，$\bm B$的特征值为$(4a-b)/6=2/3$与$b/2=0$，表是可实现的：$\bm x_1=\bm x_3$，链被完全折叠。

谱重构会毫无警告地给出这个配置，好像形状已经确定，实际上它只是无穷多个可行形状中被错误的数据选中的一个。填入其他猜测值同样会选出某个特定形状。所以缺失距离只能作为“无约束”处理，或者在明确的附加假设下另行估计，而不能伪装成观测。

\subsection{坐标重构的稳定性}
\label{sec:smap-reconstruction-stability}

距离表总带有测量误差。本小节先研究完整表上的误差怎样经过平方、双中心化和谱分解传到坐标，再说明部分距离为什么需要另作分析。所用的特征值与特征子空间扰动界取自\BookChapterRef{chap:condition-and-spectral-perturbation}{卷二第7章“条件数与谱扰动”}：对称矩阵受对称扰动$\bm E$时，每个按降序配对的特征值至多移动$\lVert\bm E\rVert_2$（Weyl界）；若前$k$个特征值与其余特征值之间有间隔$g$，且$\lVert\bm E\rVert_2\leq g/2$，前$k$个特征向量张成的子空间最大主夹角的正弦不超过$2\lVert\bm E\rVert_2/g$（Davis--Kahan界）。

\paragraph{从距离误差到Gram误差}

设观测的平方距离表为$\widetilde{\bm\Delta}=\bm\Delta+\bm E$，$\bm E$对称、零对角。双中心化是线性的，$\widetilde{\bm B}-\bm B=-\frac12\bm J\bm E\bm J$。$\bm J$是正交投影，谱范数为$1$，所以
\[
 \lVert\widetilde{\bm B}-\bm B\rVert_2\leq\frac12\lVert\bm E\rVert_2 .
\]
注意误差是加在平方距离上的：若距离本身有误差$\varepsilon_{i,j}$，平方距离的误差为$2d_{i,j}\varepsilon_{i,j}+\varepsilon_{i,j}^2$，长距离上的同样误差会被放大。

\begin{theorem}[完整距离表的谱扰动]
\label{thm:smap-mds-perturbation}
在上述记号下，设$\lambda_1\geq\cdots\geq\lambda_n$与$\widetilde\lambda_1\geq\cdots\geq\widetilde\lambda_n$分别是$\bm B$与$\widetilde{\bm B}$的特征值，则$|\widetilde\lambda_l-\lambda_l|\leq\frac12\lVert\bm E\rVert_2$。若$g=\lambda_k-\lambda_{k+1}>0$且$\lVert\bm E\rVert_2\leq g$，则两者前$k$个特征向量张成的子空间满足$\sin\theta_{\max}\leq\lVert\bm E\rVert_2/g$。
\end{theorem}
\begin{proof}
对$\bm B$与扰动$\widetilde{\bm B}-\bm B$分别应用Weyl界与Davis--Kahan界，扰动的谱范数不超过$\frac12\lVert\bm E\rVert_2$；条件$\lVert\bm E\rVert_2\leq g$保证$\frac12\lVert\bm E\rVert_2\leq g/2$，代入$2\cdot\frac12\lVert\bm E\rVert_2/g$即得。
\end{proof}

\paragraph{谱隙与小特征值}

定理分两层说明误差的去向。特征值本身稳定，误差至多为$\frac12\lVert\bm E\rVert_2$。坐标轴的方向取决于谱隙：相邻特征值接近时，很小的误差就能让坐标轴在相应的特征空间内大幅转动，这与重特征值时坐标轴本来就不唯一是同一现象的两面。

坐标的长度还要再经过开方。第$l$列坐标按$\sqrt{\lambda_l}$缩放，而
\[
 \bigl|\sqrt{\widetilde\lambda_l}-\sqrt{\lambda_l}\bigr|=\frac{|\widetilde\lambda_l-\lambda_l|}{\sqrt{\widetilde\lambda_l}+\sqrt{\lambda_l}}\leq\frac{|\widetilde\lambda_l-\lambda_l|}{\sqrt{\lambda_l}} ,
\]
特征值越小，相对误差越大。更严重的是，误差与小特征值同量级时，$\widetilde\lambda_l$可能变成负数，这个方向就整个丢失了。

\begin{example}[几乎共线的三个点]
\label{ex:smap-near-collinear}
取$\bm x_1=(-1,0)^\top$、$\bm x_2=(0,h)^\top$、$\bm x_3=(1,0)^\top$，$h=0.1$。这时$a=\delta_{1,2}=\delta_{2,3}=1.01$，$b=\delta_{1,3}=4$，由例~\ref{ex:smap-three-point-mds}的公式，特征值为$b/2=2$与$(4a-b)/6=2h^2/3\approx0.0067$。

设$d_{1,3}$被测成$2.02$，只有$1\%$的误差，于是$\widetilde\delta_{1,3}=4.0804$。第二个特征值变为$(4.04-4.0804)/6\approx-0.0067$，由正变负。事实上$2.02>2\sqrt{1.01}\approx2.0100$已违反三角不等式，表在任何维数中都不可实现。第一个特征值只从$2$变为$2.0402$，相对变化约$2\%$；纵向那个高为$0.1$的坐标则完全丢失。这里$\bm E$只有$(1,3)$、$(3,1)$两项$0.0804$，$\lVert\bm E\rVert_2=0.0804$，两个特征值的变化都在Weyl界$0.0402$之内。
\end{example}

\paragraph{比较坐标前消去合同}

评价重构误差时，不能直接比较$\widetilde{\bm Y}$与真实坐标$\bm X$的各个分量：即使距离完全精确，重构也只确定到合同为止。正确的做法是先把两组坐标都平移到质心，再求
\[
 \min_{\bm Q\in\mathrm O(d)}\lVert\widetilde{\bm Y}-\bm X\bm Q\rVert_{\mathrm F},
\]
它的解由第~\ref{sec:smap-rotation}节的正交对齐给出。这里必须允许反射，因为距离同样无法辨认反射。

\paragraph{部分距离的稳定性}

只观测部分距离时，没有完整的$\bm B$可以分析，稳定性取决于观测约束本身是否足够，不能由某张完整表的谱隙推出。仍以两个共边三角形为例，把$\bm x_4$改为$(1,-\eta)^\top$，$\eta>0$很小，第四个点贴近公共边。由$d_{1,4}^2=1+\eta^2$，若$d_{1,4}^2$与$d_{2,4}^2$都有误差$\varepsilon$，$\eta^2$就改变$\varepsilon$，$\eta$约改变$\varepsilon/(2\eta)$，放大倍数随$\eta\to0$无界增长。

同时，$\bm x_4$在公共边两侧的两个解只相距$2\eta$，误差稍大就无法判断它在哪一侧。这些现象只与观测图在这组点上的约束几何有关，而与真实配置的完整距离表是否有良好的谱隙无关：真实配置的完整表在这个例子中完全正常。

\section{空间映射的估计}
\label{sec:smap-map-estimation}

现在转向第二类观测：同一组点在两套坐标中都有记录，要求出两套坐标之间的映射。第一步是选定允许的映射类型，因为有限个点对本身几乎不约束任意的映射；选定以后，每一对对应点都给参数加上若干条件，精确数据下是方程，含噪数据下是拟合残差。本节依次处理刚体、仿射和射影三种模型。它们保留的几何结构不同，参数的约束、求解方法与合适的误差也随之不同。

\subsection{对应关系与映射模型}
\label{sec:smap-correspondence-models}

设源点$\bm x_1,\ldots,\bm x_n\in\mathbb R^d$与目标点$\bm y_1,\ldots,\bm y_n\in\mathbb R^m$按下标一一对应。一对对应点给未知映射$f:\mathbb R^d\to\mathbb R^m$加上一个条件$f(\bm x_i)=\bm y_i$，即$m$个标量方程。$n$对点共给出$nm$个方程，这个数是有限的，而一般映射的自由度无穷多。

\paragraph{有限对应不能确定任意映射}

即使只要求光滑，有限个点对也确定不了映射。设$f$满足全部条件，取任意向量$\bm c\in\mathbb R^m$，令$\varphi(\bm x)=\prod_{i=1}^n\lVert\bm x-\bm x_i\rVert_2^2$，则$f+\varphi\,\bm c$光滑，并且在每个$\bm x_i$处与$f$取值相同。因此必须事先指定映射族$\mathcal F$，只在$\mathcal F$中寻找$f$。所选的族同时规定了定义域、值域和映射必须保留的结构。

\paragraph{三种映射模型}

本节使用表~\ref{tab:smap-models}中的三种模型。参数个数说明至少需要多少对应点：每对点提供$m$个方程，刚体与仿射模型中的方程关于参数分别是带正交约束的和线性的，射影模型中的方程则只在乘以一个非零倍数后成立。

\begin{table}[htbp]
 \centering
 \caption{三种映射模型。刚体模型要求$m=d$，射影模型这里只讨论平面到平面的情形，$\widetilde{\bm x}$为齐次坐标，见第~\ref{sec:smap-projective}节。}
 \label{tab:smap-models}
 \begin{tabular}{@{}llll@{}}
 \toprule
 模型 & 形式 & 参数个数 & 保留的结构\\
 \midrule
 刚体 & $\bm R\bm x+\bm t$，$\bm R\in\mathrm{SO}(d)$ & $d(d+1)/2$ & 距离与定向\\
 仿射 & $\bm A\bm x+\bm b$，$\bm A\in\mathbb R^{m\times d}$ & $m(d+1)$ & 平行与共线点的比例\\
 射影 & $[\widetilde{\bm x}]\mapsto[\bm H\widetilde{\bm x}]$，$\bm H$可逆 & $8$ & 共线\\
 \bottomrule
 \end{tabular}
\end{table}

\paragraph{可满足、唯一与稳定}

选定$\mathcal F$后，有三个需要分别回答的问题。第一，数据是否可满足：是否存在$f\in\mathcal F$使$f(\bm x_i)=\bm y_i$对所有$i$精确成立；含噪数据一般不可满足，这时改为最小化残差。第二，参数是否唯一：精确可满足时解是否只有一个，最小化残差时极小点是否只有一个。第三，结果是否稳定：数据有小扰动时，所得映射是否只有小变化。三个问题互不蕴含，后面的每种模型都要分别核对。

对应关系在本节中始终是给定的输入。若对应本身有误，例如把第三个点错配给第五个点，所有结论都针对错误的数据成立，估计方法本身不会发现这种错误。

\subsection{刚体变换估计}
\label{sec:smap-rigid}

刚体模型把源点送到$\bm R\bm x_i+\bm t$，其中$\bm R\in\mathrm{SO}(d)$，$\bm t\in\mathbb R^d$，这就是第~\ref{chap:lie-groups}章的刚体位姿。给定两组按下标对应的欧氏点，以平方距离残差
\begin{equation}
 F(\bm R,\bm t)=\sum_{i=1}^n\lVert\bm y_i-\bm R\bm x_i-\bm t\rVert_2^2
 \label{eq:smap-rigid-objective}
\end{equation}
衡量拟合程度。未知量$\bm R$受正交约束，不能像普通最小二乘那样直接求导解线性方程。下面先消去平移，再把旋转问题化为一个矩阵上的迹最大化，最后讨论数据何时足以确定变换。

\subsubsection{平移的消除}
\label{sec:smap-translation}

先固定$\bm R$，看最优平移是什么。记$\bm a_i=\bm y_i-\bm R\bm x_i$，目标变为$\sum_i\lVert\bm a_i-\bm t\rVert_2^2$，即找一个点使它到$\bm a_1,\ldots,\bm a_n$的平方距离和最小。

\begin{lemma}[最优平移]
\label{lem:smap-optimal-translation}
记$\overline{\bm x}=\frac1n\sum_i\bm x_i$，$\overline{\bm y}=\frac1n\sum_i\bm y_i$，$\bm x_i'=\bm x_i-\overline{\bm x}$，$\bm y_i'=\bm y_i-\overline{\bm y}$。对固定的$\bm R$，$F(\bm R,\cdot)$的唯一极小点是$\bm t^\star(\bm R)=\overline{\bm y}-\bm R\overline{\bm x}$，并且
\[
 F(\bm R,\bm t^\star(\bm R))=\sum_{i=1}^n\lVert\bm y_i'-\bm R\bm x_i'\rVert_2^2 .
\]
\end{lemma}
\begin{proof}
记$\overline{\bm a}=\frac1n\sum_i\bm a_i=\overline{\bm y}-\bm R\overline{\bm x}$。由$\sum_i(\bm a_i-\overline{\bm a})=\bm0$，交叉项求和为零，
\[
 \sum_i\lVert\bm a_i-\bm t\rVert_2^2=\sum_i\lVert\bm a_i-\overline{\bm a}\rVert_2^2+n\lVert\overline{\bm a}-\bm t\rVert_2^2 ,
\]
所以唯一极小点是$\bm t=\overline{\bm a}$。代入后$\bm a_i-\overline{\bm a}=\bm y_i'-\bm R\bm x_i'$。
\end{proof}

引理把原问题化为中心化点之间的纯旋转拟合。中心化消去的只是平移这个未知量；点与点之间的配对仍然是输入，质心本身与配对无关，但下一小节的交叉矩阵逐对相乘，完全依赖于配对。本节不讨论配对的求法。

\subsubsection{旋转的估计}
\label{sec:smap-rotation}

展开中心化后的目标。旋转保持长度，所以$\lVert\bm R\bm x_i'\rVert_2=\lVert\bm x_i'\rVert_2$，
\[
 \sum_i\lVert\bm y_i'-\bm R\bm x_i'\rVert_2^2=\sum_i\lVert\bm y_i'\rVert_2^2+\sum_i\lVert\bm x_i'\rVert_2^2-2\sum_i\bm y_i'^\top\bm R\bm x_i' .
\]
前两项与$\bm R$无关。最后一项可以用迹写成$\sum_i\bm y_i'^\top\bm R\bm x_i'=\operatorname{tr}(\bm R^\top\bm M)$，其中
\begin{equation}
 \bm M=\sum_{i=1}^n\bm y_i'\bm x_i'^\top\in\mathbb R^{d\times d}
 \label{eq:smap-cross-matrix}
\end{equation}
称为两组中心化点的交叉矩阵。于是旋转估计等价于在$\mathrm{SO}(d)$中最大化$\operatorname{tr}(\bm R^\top\bm M)$。

这个问题已在卷二的矩阵逼近一章解决：由$\lVert\bm M-\bm R\rVert_{\mathrm F}^2=\lVert\bm M\rVert_{\mathrm F}^2+d-2\operatorname{tr}(\bm R^\top\bm M)$，最大化迹就是找离$\bm M$最近的正交矩阵或旋转。那里的结论是：设$\bm M=\bm U\bm\Sigma\bm V^\top$为奇异值分解，在全部正交矩阵中，$\operatorname{tr}(\bm R^\top\bm M)\leq\sum_i\sigma_i$，由$\bm U\bm V^\top$取到，它就是$\bm M$的极分解的正交因子；只允许旋转时，记$s=\det(\bm U\bm V^\top)$，最优解为$\bm U\operatorname{diag}(1,\ldots,1,s)\bm V^\top$，迹的最大值为$\sigma_1+\cdots+\sigma_{d-1}+s\sigma_d$。极分解本身可参考\BookChapterRef{chap:polar-decomposition}{卷二第6章“极分解”}。

\begin{theorem}[对应点的刚体估计]
\label{thm:smap-rigid-estimate}
设$\bm M=\bm U\bm\Sigma\bm V^\top$，$s=\det(\bm U\bm V^\top)$。则
\[
 \bm R^\star=\bm U\operatorname{diag}(1,\ldots,1,s)\bm V^\top,\qquad\bm t^\star=\overline{\bm y}-\bm R^\star\overline{\bm x}
\]
使式~\eqref{eq:smap-rigid-objective}在$\mathrm{SO}(d)\times\mathbb R^d$上取最小值，最小值为
\[
 \sum_i\lVert\bm x_i'\rVert_2^2+\sum_i\lVert\bm y_i'\rVert_2^2-2\Bigl(\sum_{l=1}^{d-1}\sigma_l+s\,\sigma_d\Bigr).
\]
若把$\mathrm{SO}(d)$换成$\mathrm O(d)$，则取$\bm R^\star=\bm U\bm V^\top$，最小值中的括号换成$\sum_l\sigma_l$。
\end{theorem}
\begin{proof}
由引理~\ref{lem:smap-optimal-translation}，对每个$\bm R$，最优平移是$\bm t^\star(\bm R)$，剩下的目标是上面的展开式。它关于$\bm R$的最小化等价于$\operatorname{tr}(\bm R^\top\bm M)$的最大化，由上面引述的结论即得最优旋转与最大迹，代入展开式得最小值。正交情形同理。
\end{proof}

计算步骤是：求两组点的质心并中心化；按式~\eqref{eq:smap-cross-matrix}累加交叉矩阵；求$\bm M$的奇异值分解并按需要作行列式修正；恢复平移；最后把$\bm R^\star\bm x_i+\bm t^\star$与$\bm y_i$逐个比较，报告残差。

\begin{example}[平面中的旋转与平移]
\label{ex:smap-rigid-rotation}
取$\bm x_1=(1,0)^\top$、$\bm x_2=(0,2)^\top$、$\bm x_3=(-1,0)^\top$、$\bm x_4=(0,-2)^\top$，质心为零。目标点由逆时针旋转$90^\circ$再平移$(3,1)^\top$得到：$\bm y_1=(3,2)^\top$、$\bm y_2=(1,1)^\top$、$\bm y_3=(3,0)^\top$、$\bm y_4=(5,1)^\top$，质心为$(3,1)^\top$，中心化后为$(0,1)^\top$、$(-2,0)^\top$、$(0,-1)^\top$、$(2,0)^\top$。

逐项累加得$\bm M=\left[\begin{smallmatrix}0&-8\\2&0\end{smallmatrix}\right]$。它可以写成$\bm M=\bm R_{90}\operatorname{diag}(2,8)$，其中$\bm R_{90}=\left[\begin{smallmatrix}0&-1\\1&0\end{smallmatrix}\right]$，这正是极分解：对称因子$\operatorname{diag}(2,8)$正定，正交因子$\bm R_{90}$的行列式为$1$，所以$\bm R^\star=\bm R_{90}$，无需修正。平移$\bm t^\star=(3,1)^\top-\bm R_{90}\bm0=(3,1)^\top$，残差为零，恢复了生成数据的变换。
\end{example}

\paragraph{正交对齐与旋转对齐}

两种约束用于不同的场合。比较第~\ref{sec:smap-distance-reconstruction}节的两组重构坐标时，距离本来就辨认不了反射，应在$\mathrm O(d)$中对齐，否则会把一个反射误算成误差。估计物体的刚体位姿时，反射不是物理上可实现的运动，应在$\mathrm{SO}(d)$中对齐。数据含反射时，两者的结果截然不同。

\begin{example}[镜像数据的两种对齐]
\label{ex:smap-mirror}
沿用上例的$\bm x_i$，令$\bm y_i$为$\bm x_i$关于横轴的镜像：$(1,0)^\top$、$(0,-2)^\top$、$(-1,0)^\top$、$(0,2)^\top$。交叉矩阵为$\bm M=\operatorname{diag}(2,-8)$，奇异值为$8$与$2$，$\det\bm M<0$。在$\mathrm O(2)$中，最优解是反射$\operatorname{diag}(1,-1)$，残差为零。

在$\mathrm{SO}(2)$中，$s=-1$，最大迹为$8-2=6$。直接验证：旋转$\bm R(\theta)$给出$\operatorname{tr}(\bm R^\top\bm M)=2\cos\theta-8\cos\theta=-6\cos\theta$，在$\theta=\pi$处取最大值$6$，即$\bm R^\star=-\bm I$。最小残差为$10+10-2\cdot6=8$。一个不允许反射的模型无法解释镜像数据，残差如实反映了这一点。
\end{example}

\subsubsection{退化配置}
\label{sec:smap-degenerate}

刚体估计总能算出一个$\bm R^\star$，但它未必是唯一的。先看三维中的两种极端配置。

若全部源点重合，中心化后$\bm x_i'=\bm0$，$\bm M=\bm0$，每个旋转都使迹为零，都是最优解；平移$\bm t^\star=\overline{\bm y}-\bm R\overline{\bm x}$随所选的$\bm R$而变。数据只确定了那一个点被送到哪里，绕它的任何转动都无法辨认。

若源点共线，写$\bm x_i'=a_i\bm u$，$\bm u$为单位向量，则$\bm M=\bm w\bm u^\top$，$\bm w=\sum_ia_i\bm y_i'$，秩至多为$1$。这时$\operatorname{tr}(\bm R^\top\bm M)=(\bm R\bm u)^\top\bm w$只依赖$\bm R\bm u$；$\bm w\neq\bm0$时，最优解是全部满足$\bm R\bm u=\bm w/\lVert\bm w\rVert_2$的旋转，它们还可以绕$\bm w$方向任意转动，构成一个单参数族。几何上这很直观：绕点所在直线的转动不移动任何一个点。

\begin{theorem}[旋转估计的唯一性]
\label{thm:smap-rotation-uniqueness}
设$\bm M=\bm U\bm\Sigma\bm V^\top$，$\sigma_1\geq\cdots\geq\sigma_d\geq0$，$s=\det(\bm U\bm V^\top)$。
\begin{enumerate}
 \item 若$\operatorname{rank}\bm M\leq d-2$，则$\mathrm{SO}(d)$中的最优旋转不唯一。
 \item 若$\bm M$可逆且$s=1$，则最优旋转唯一，为$\bm U\bm V^\top$。
\end{enumerate}
\end{theorem}
\begin{proof}
(1) 此时$\sigma_{d-1}=\sigma_d=0$，$\bm\Sigma$的最后两行为零。设$\bm R^\star$最优，对任意$\bm Z_0\in\mathrm{SO}(2)$，令$\bm Z=\bm U\operatorname{diag}(\bm I_{d-2},\bm Z_0)\bm U^\top\in\mathrm{SO}(d)$。由于$\operatorname{diag}(\bm I_{d-2},\bm Z_0^\top)\bm\Sigma=\bm\Sigma$，
\[
 \operatorname{tr}\bigl((\bm Z\bm R^\star)^\top\bm M\bigr)=\operatorname{tr}\bigl(\bm R^{\star\top}\bm U\operatorname{diag}(\bm I_{d-2},\bm Z_0^\top)\bm\Sigma\bm V^\top\bigr)=\operatorname{tr}(\bm R^{\star\top}\bm M),
\]
所以$\bm Z\bm R^\star$也最优；$\bm Z_0\neq\bm I$时$\bm Z\neq\bm I$，这给出无穷多个不同的最优解。

(2) 对任一旋转$\bm R$，令$\bm Z=\bm U^\top\bm R\bm V$，它正交，于是$\operatorname{tr}(\bm R^\top\bm M)=\operatorname{tr}(\bm Z^\top\bm\Sigma)=\sum_i\sigma_iz_{i,i}\leq\sum_i\sigma_i$，因为单位列向量的分量$|z_{i,i}|\leq1$。全部$\sigma_i>0$，等号要求每个$z_{i,i}=1$，单位列向量$\bm Z_{:,i}$只能是$\bm e_i$，所以$\bm Z=\bm I$，$\bm R=\bm U\bm V^\top$；$s=1$保证它是旋转。
\end{proof}

定理没有覆盖的情形也有明确的答案：一般地，最优旋转唯一当且仅当$\operatorname{rank}\bm M\geq d-1$，并且在$s=-1$时还要求$\sigma_{d-1}>\sigma_d$。这里不给出证明。对精确的刚体数据，$\bm y_i'=\bm R\bm x_i'$，$\bm M=\bm R\sum_i\bm x_i'\bm x_i'^\top$，秩等于中心化源点张成的维数，并且$s=1$。所以在三维中，三个不共线的点已足以唯一确定旋转，共面并不构成退化。

\paragraph{数据退化、残差与参数冗余}

三种现象需要分开。数据退化是指$\bm M$的秩不足，最优解成族，这由点的几何配置决定，与噪声无关。含噪残差是最优值本身不为零，它度量模型与数据的不符，与解是否唯一无关。坐标表示冗余是指同一个旋转有多种参数，例如单位四元数$\bm q$与$-\bm q$表示同一个旋转，变换唯一而参数不唯一。

较小的拟合残差不能解释为变换已被唯一确定。共线的精确数据对整整一族旋转都给出零残差。接近退化时，例如三维中几乎共线的点，$\sigma_{d-1}$很小，最优旋转虽然唯一，却会随数据的微小扰动绕点所在直线大幅转动，这与第~\ref{sec:smap-reconstruction-stability}节中谱隙的作用是同一回事。

\subsection{仿射映射估计}
\label{sec:smap-affine}

刚体模型只允许转动和平移。若两套坐标之间还有伸缩或剪切，例如两种测绘坐标系的单位不同、坐标轴不垂直，就需要更大的模型族。把对应关系写为
\[
 \bm y_i\approx\bm A\bm x_i+\bm b,\qquad\bm A\in\mathbb R^{m\times d},\ \bm b\in\mathbb R^m .
\]
$\bm A$没有任何约束，可以伸缩、剪切，也可以降维，例如$m=2<d=3$时的平行投影。只有当$m=d$且明确要求两套坐标之间可逆换算时，才另加$\bm A$可逆的条件。

\subsubsection{增广坐标求解}
\label{sec:smap-augmented}

仿射映射关于$(\bm A,\bm b)$是线性的，只差一个常数项。在每个源点坐标后补上常数$1$，记$\widetilde{\bm x}_i=(\bm x_i^\top,1)^\top\in\mathbb R^{d+1}$，称为\term{增广坐标}（augmented coordinates），再把参数合并为$\bm W=[\bm A\ \bm b]\in\mathbb R^{m\times(d+1)}$，就有$\bm A\bm x_i+\bm b=\bm W\widetilde{\bm x}_i$。

把$\widetilde{\bm x}_i^\top$按行排成设计矩阵$\widetilde{\bm X}\in\mathbb R^{n\times(d+1)}$，$\bm y_i^\top$按行排成$\bm Y\in\mathbb R^{n\times m}$。欧氏输出残差为
\[
 \sum_{i=1}^n\lVert\bm y_i-\bm W\widetilde{\bm x}_i\rVert_2^2=\lVert\bm Y-\widetilde{\bm X}\bm W^\top\rVert_{\mathrm F}^2=\sum_{k=1}^m\lVert\bm Y_{:,k}-\widetilde{\bm X}\bm w_k\rVert_2^2 ,
\]
其中$\bm w_k^\top$是$\bm W$的第$k$行。目标分解为$m$个共用同一设计矩阵的普通最小二乘问题，每个目标坐标单独求解，最小二乘解满足正规方程$\widetilde{\bm X}^\top\widetilde{\bm X}\bm W^\top=\widetilde{\bm X}^\top\bm Y$。最小二乘解的构造、正规方程与解的核对可参考\BookChapterRef{chap:linear-solution-selection}{卷二第9章“线性方程与最小二乘的选解”}。

\paragraph{精确可解与含噪拟合}

精确可解要求$\bm Y$的每一列都落在$\widetilde{\bm X}$的列空间中。平面上取三个不共线的点，$\widetilde{\bm X}$是可逆的$3\times3$矩阵，任意给定的三个目标点都能被唯一的仿射映射精确满足。点数多于参数所需的最少数量时，任意给定的目标坐标一般不再能被同一个仿射映射精确满足，因为仿射映射保持点之间的仿射关系。

\begin{example}[正方形四顶点的仿射拟合]
\label{ex:smap-affine-square}
取源点$(0,0)$、$(1,0)$、$(0,1)$、$(1,1)$，它们满足$\bm x_1+\bm x_4=\bm x_2+\bm x_3$。任何仿射映射都保持这一关系，要求$\bm y_1+\bm y_4=\bm y_2+\bm y_3$。取目标点$(0,0)$、$(1,0)$、$(0,1)$、$(2,2)$，左右相差$(1,1)$，所以不可精确满足。

记$\bm c=(1,-1,-1,1)^\top$，它与$\widetilde{\bm X}$的三列都正交，而$\widetilde{\bm X}$的列空间是三维的，所以$\bm c$张成列空间的正交补。最小二乘残差是$\bm Y$在$\bm c$上的投影，$\bm c^\top\bm Y=(1,1)$，残差矩阵为$\frac14\bm c\,(1,1)$，每个点的残差为$\pm(1/4,1/4)$，残差平方和为$1/2$。

拟合值为$(-1/4,-1/4)$、$(5/4,1/4)$、$(1/4,5/4)$、$(7/4,7/4)$，由此读出$\bm b=(-1/4,-1/4)^\top$，$\bm A$的两列分别为$(3/2,1/2)^\top$与$(1/2,3/2)^\top$。代回检查$\bm b+\bm A(1,1)^\top=(7/4,7/4)^\top$，与第四个拟合值一致。不相容的量被平均地摊给四个点，而不是全部归到第四个点上。
\end{example}

\subsubsection{唯一性与可逆性}
\label{sec:smap-affine-uniqueness}

最小二乘参数唯一，当且仅当设计矩阵$\widetilde{\bm X}$满列秩。下面的定理把这个代数条件翻译成源点的几何配置。

\begin{theorem}[增广设计矩阵的满列秩]
\label{thm:smap-affine-span}
$\widetilde{\bm X}$满列秩，当且仅当差向量$\bm x_2-\bm x_1,\ldots,\bm x_n-\bm x_1$张成$\mathbb R^d$，即源点仿射张成整个定义域。
\end{theorem}
\begin{proof}
$\widetilde{\bm X}(\bm a;\beta)=\bm0$的意思是$\bm a^\top\bm x_i+\beta=0$对所有$i$成立，相减得$\bm a^\top(\bm x_i-\bm x_1)=0$。若差向量张成$\mathbb R^d$，则$\bm a=\bm0$，进而$\beta=0$，核为零。反之，若差向量不张成$\mathbb R^d$，取与它们都正交的非零$\bm a$，令$\beta=-\bm a^\top\bm x_1$，就得到核中的非零向量。
\end{proof}

于是平面上至少需要三个不共线的点，空间中至少需要四个不共面的点。条件满足时参数唯一，但未必稳定：$\widetilde{\bm X}$的最小奇异值很小时，数据的微小误差会被大幅放大。例如取源点$(0,0)$、$(1,0)$、$(0,\varepsilon)$，精确拟合给出$\bm A$的第二列$(\bm y_3-\bm y_1)/\varepsilon$，第三个目标点的误差$\eta$使这一列改变$\eta/\varepsilon$。源点几乎共线时，垂直于那条直线方向上的伸缩几乎无法测定。

\paragraph{映射的可逆性另需核验}

即使参数唯一，所得的映射也未必可逆。取源点$(0,0)$、$(1,0)$、$(0,1)$，目标点$(0,0)$、$(1,0)$、$(2,0)$，精确解为$\bm b=\bm0$、$\bm A=\left[\begin{smallmatrix}1&2\\0&0\end{smallmatrix}\right]$，它由数据唯一确定，却把整个平面压到横轴上，$\det\bm A=0$。

所以“观测足以确定映射”与“所确定的映射保留维数”是两个问题。前者看源点的配置，即$\widetilde{\bm X}$的秩与条件数；后者看估计出的$\bm A$本身，例如它的行列式或最小奇异值。应用要求同维可逆换算时，必须在估计之后单独核验$\bm A$，必要时报告它接近奇异的程度。

\subsection{射影变换估计}
\label{sec:smap-projective}

从两个不同位置拍摄同一张平面海报，两张照片中对应点之间的关系既不是刚体的，也不是仿射的：海报上的平行线在照片中可能汇聚，等距的点在照片中可能疏密不均。这种关系由射影变换描述。改用第~\ref{chap:lie-groups}章的齐次坐标，平面点$(u,v)$记为$\widetilde{\bm x}=(u,v,1)^\top$，凡与它成非零倍数的向量都代表同一个射影点。

射影变换由可逆矩阵$\bm H\in\mathbb R^{3\times3}$诱导，把射影点$[\widetilde{\bm x}]$送到$[\bm H\widetilde{\bm x}]$。$\bm H$与$c\bm H$（$c\neq0$）给出同一个变换，所以九个元素中只有八个自由度。这类变换只保持共线关系，不保持距离、比例和平行，刚体模型的正交约束无从谈起，必须另建约束。还要区分两种误差：齐次方程上的代数残差，以及还原成平面坐标后的重投影误差。

\subsubsection{齐次约束与可辨识性}
\label{sec:smap-homogeneous-constraints}

对应点要求$\bm H\widetilde{\bm x}_i$与$\widetilde{\bm y}_i=(u_i',v_i',1)^\top$成比例。记$\bm H$的三行为$\bm h_1^\top,\bm h_2^\top,\bm h_3^\top$，比例关系等价于两个方程
\begin{equation}
 \bm h_1^\top\widetilde{\bm x}_i-u_i'\,\bm h_3^\top\widetilde{\bm x}_i=0,\qquad
 \bm h_2^\top\widetilde{\bm x}_i-v_i'\,\bm h_3^\top\widetilde{\bm x}_i=0 .
 \label{eq:smap-dlt-rows}
\end{equation}
确实，若$\bm h_3^\top\widetilde{\bm x}_i=c\neq0$，两式说明$\bm H\widetilde{\bm x}_i=c\,\widetilde{\bm y}_i$；若$c=0$，两式迫使$\bm H\widetilde{\bm x}_i=\bm0$，可逆的$\bm H$不会如此。

式~\eqref{eq:smap-dlt-rows}关于$\bm h=(\bm h_1;\bm h_2;\bm h_3)\in\mathbb R^9$是线性齐次的。每对点给出两行，$n$对点叠成$\bm A\bm h=\bm0$，$\bm A\in\mathbb R^{2n\times9}$。这种由对应点直接建立线性方程的做法称为\term{直接线性变换}（direct linear transformation）。齐次方程总有零解，有意义的问题是：$\bm A$的核是否恰好是一维的，并且核中的方向是否对应可逆矩阵。四对点给出八个方程，是核为一维所需的最少数量。

\begin{theorem}[四对点确定射影变换]
\label{thm:smap-four-points}
设平面上的源点$\bm x_1,\ldots,\bm x_4$中任意三点不共线，目标点$\bm y_1,\ldots,\bm y_4$中任意三点也不共线。则式~\eqref{eq:smap-dlt-rows}构成的$8\times9$矩阵$\bm A$的核是一维的，核中的非零向量对应可逆矩阵$\bm H$，并且$\bm H$是唯一（差一个非零倍数）把$[\widetilde{\bm x}_i]$送到$[\widetilde{\bm y}_i]$的射影变换。
\end{theorem}
\begin{proof}
三点不共线等价于它们的齐次坐标线性无关，所以$\widetilde{\bm x}_1,\widetilde{\bm x}_2,\widetilde{\bm x}_3$构成$\mathbb R^3$的基。写$\widetilde{\bm x}_4=\sum_{j\leq3}\alpha_j\widetilde{\bm x}_j$，每个$\alpha_j\neq0$：例如$\alpha_3=0$意味着$\widetilde{\bm x}_4$在$\widetilde{\bm x}_1,\widetilde{\bm x}_2$的张成中，即$\bm x_1,\bm x_2,\bm x_4$共线。令$\bm P=[\alpha_1\widetilde{\bm x}_1,\alpha_2\widetilde{\bm x}_2,\alpha_3\widetilde{\bm x}_3]$，它可逆，把$\bm e_j$送到$\alpha_j\widetilde{\bm x}_j$，把$(1,1,1)^\top$送到$\widetilde{\bm x}_4$。对目标点同样构造$\bm Q$，系数记为$\beta_j$。

设$\bm h$在$\bm A$的核中，对应矩阵$\bm H$。由式~\eqref{eq:smap-dlt-rows}后的说明，$\bm H\widetilde{\bm x}_i=\mu_i\widetilde{\bm y}_i$，$\mu_i$可以为零。令$\bm K=\bm Q^{-1}\bm H\bm P$，则$\bm K\bm e_j=\alpha_j\mu_j\bm Q^{-1}\widetilde{\bm y}_j=(\alpha_j\mu_j/\beta_j)\bm e_j$，所以$\bm K$是对角矩阵；又$\bm K(1,1,1)^\top=\mu_4\bm Q^{-1}\widetilde{\bm y}_4=\mu_4(1,1,1)^\top$，三个对角元都等于$\mu_4$。于是$\bm H=\mu_4\bm Q\bm P^{-1}$，核由$\bm Q\bm P^{-1}$张成，是一维的，非零元对应的矩阵可逆。任何满足对应的射影变换的矩阵都在核中，唯一性随之成立。
\end{proof}

\begin{example}[正方形到一般四边形]
\label{ex:smap-homography-square}
把单位正方形的顶点$(0,0)$、$(1,0)$、$(0,1)$、$(1,1)$送到$(0,0)$、$(1,0)$、$(0,1)$、$(2,2)$，两组点中都没有三点共线。可以验证
\[
 \bm H=\begin{bmatrix}2&0&0\\0&2&0\\-1&-1&3\end{bmatrix}
\]
满足要求：它把$(1,0,1)^\top$送到$(2,0,2)^\top$，即$(1,0)$；把$(1,1,1)^\top$送到$(2,2,1)^\top$，即$(2,2)$；另两点同样可查。$\det\bm H=12\neq0$。由定理~\ref{thm:smap-four-points}，它在差一个倍数的意义下是唯一解。上一小节的仿射模型无法拟合这组数据，射影模型却能精确满足。
\end{example}

\paragraph{退化配置与不合法的解}

点的配置违反定理条件时，核仍然非零，却可能不含合法的射影变换。设源点$\bm x_1,\bm x_2,\bm x_3$共线，目标点中却没有三点共线。可逆的$\bm H$把共线点送到共线点，所以不存在可逆解；但八个方程、九个未知量的齐次方程总有非零解，这些解对应的$\bm H$都是奇异的，会把某些点送到零向量，或把平面压到一条直线上。

所以求得齐次解以后，必须核验$\bm H$是否可逆，不能把任意非零齐次解都当作合法的射影变换。点数多于四对时，也要检查是否大多数点落在同一条直线附近，这时虽然方程很多，独立的约束却不够。

\subsubsection{归一化代数求解}
\label{sec:smap-normalized-dlt}

点数多于四对且含噪时，$\bm A\bm h=\bm0$一般只有零解，需要改为拟合。直接最小化$\lVert\bm A\bm h\rVert_2$总得到$\bm h=\bm0$，必须先固定尺度。固定$h_{3,3}=1$会排除真实变换恰好有$h_{3,3}=0$的情形，即源点原点被送到无穷远；固定$\lVert\bm h\rVert_2=1$则对各分量一视同仁。

后一种选择就是齐次最小二乘：在单位向量中最小化$\lVert\bm A\bm h\rVert_2^2$，最优解是$\bm A^\top\bm A$最小特征值的特征方向，也就是$\bm A$的最小奇异值对应的右奇异向量；最小特征值为单重时，解只差一个符号。这一结论在前面引用的卷二线性方程一章中已经证明。

\paragraph{坐标归一化}

直接用原始坐标建立$\bm A$，数值上往往很差。设像素坐标的量级为$10^3$，则$\bm A$的一行中既有常数$1$，又有$u'u$这样量级为$10^6$的项，各列尺度相差悬殊，最小奇异值方向对误差极为敏感。因此先对两组点分别作相似变换$\bm T$与$\bm T'$：平移使质心到原点，再伸缩使各点到原点的平均距离为$\sqrt2$。

完整的步骤是：计算$\widehat{\bm x}_i=\bm T\widetilde{\bm x}_i$、$\widehat{\bm y}_i=\bm T'\widetilde{\bm y}_i$；用它们建立$\widehat{\bm A}$并求单位范数的齐次最小二乘解$\widehat{\bm H}$；还原$\bm H=\bm T'^{-1}\widehat{\bm H}\bm T$；最后检查$\bm H$是否可逆，例如比较它的最小奇异值与最大奇异值。

\paragraph{归一化改变了代数目标}

归一化不只是数值技巧，它改变了所求的问题。精确数据下二者给出同一个变换：若$\bm H\widetilde{\bm x}_i\propto\widetilde{\bm y}_i$，则$(\bm T'\bm H\bm T^{-1})\widehat{\bm x}_i\propto\widehat{\bm y}_i$，核与坐标系的选择无关。含噪数据下，$\lVert\widehat{\bm A}\widehat{\bm h}\rVert_2$一般不是$\lVert\bm A\bm h\rVert_2$的常数倍，两个目标的最优解也就不同。

代数残差的大小依赖于坐标的表示方式，而不是一个几何上固定的量。归一化选择了一种数值条件较好、对平移和伸缩不敏感的表示，但“代数最优”始终是相对于这一表示而言的。要得到不依赖表示的评价，需要下一小节的几何误差。

\subsubsection{重投影误差}
\label{sec:smap-reprojection}

射影点要还原成平面坐标，需要选定仿射图：把$\bm z=(z_1,z_2,z_3)^\top$送到$\pi(\bm z)=(z_1/z_3,z_2/z_3)$，只在$z_3\neq0$时有定义。源点$\bm x_i$的映射点为$\widehat{\bm y}_i=\pi(\bm H\widetilde{\bm x}_i)$，先要核验分母$c_i=\bm h_3^\top\widetilde{\bm x}_i\neq0$。在例~\ref{ex:smap-homography-square}中，$c=-u-v+3$，直线$u+v=3$上的点被送到无穷远，例如$(3/2,3/2)$没有平面上的像。

分母非零时，定义\term{重投影误差}（reprojection error）
\[
 E(\bm H)=\sum_{i=1}^n\lVert\bm y_i-\pi(\bm H\widetilde{\bm x}_i)\rVert_2^2 ,
\]
它是目标平面上的欧氏距离之和，与$\bm H$乘以非零常数无关，也与齐次坐标的书写方式无关。

\paragraph{代数残差与几何误差的关系}

两种误差之间有精确的关系。式~\eqref{eq:smap-dlt-rows}第一式的左端为$\bm h_1^\top\widetilde{\bm x}_i-u_i'c_i=c_i\bigl(\pi_1(\bm H\widetilde{\bm x}_i)-u_i'\bigr)$，第二式同理，所以第$i$对点的两行代数残差恰为$c_i(\widehat{\bm y}_i-\bm y_i)$，从而
\[
 \lVert\bm A\bm h\rVert_2^2=\sum_{i=1}^nc_i^2\,\lVert\widehat{\bm y}_i-\bm y_i\rVert_2^2 .
\]
\term{代数残差}（algebraic residual）是以$c_i^2$加权的重投影误差。权重随点和$\bm H$而变：接近被送到无穷远的点，$c_i$小，几何误差再大，在代数目标中也几乎不被计较。

所以代数最优解一般不使重投影误差最小，二者只在精确数据下同时为零。实际的做法是先用归一化的直接线性变换得到初值，再以重投影误差为目标迭代改进；这是一个非线性最小二乘问题，其方法与收敛分析可参考\BookChapterRef{chap:smooth-optimization}{卷五第3章“光滑优化”}。两组点都含噪时，还可以同时考虑正向与反向的重投影误差，选择哪种几何误差取决于噪声出现在哪一端。

\section{本章小结}
\label{sec:smap-summary}

本章开头问：只知道相对距离，能否安排坐标；知道两组对应点，能否求出它们之间的映射。两个问题的第一步相同，都是辨认观测不能确定的东西。距离对整体平移、旋转和反射一无所知，因此最多把点配置确定到合同为止；缺失的距离还会留下真正的形变，例如只测边长的正方形可以压成菱形。固定参照能消去合同的自由度，却补不出缺失的观测；用仪器测得的锚点则是额外的观测，不是约定。

完整距离表的重构由双中心化完成。平方距离经过$-\frac12\bm J\bm\Delta\bm J$变成以质心为原点的Gram矩阵，它半正定当且仅当距离表欧氏可实现，秩就是所需的最小维数，谱分解直接给出坐标。完整距离相同的两组点必定合同，所以合同是完整表下唯一的不确定性。指定低维时，保留最大的正特征值在Gram矩阵的Frobenius误差下最优，负谱只会增加误差；但距离误差是另一个量，应力准则下截断坐标一般不是最优的。

部分距离的问题由观测图描述。连通不保证形状确定，局部刚性也不保证全局唯一，两个共边三角形就有互为反射的两个实现。缺边、不相容与含噪是三种不同的困难，残差只能建在已观测的边上；把缺失填成零等于断言两点重合，会让谱重构给出一个看似确定的错误形状。稳定性方面，完整表的误差经双中心化至多减半地传到Gram矩阵，特征值稳定，坐标轴方向取决于谱隙，小特征值经开方后最为脆弱；部分距离的稳定性则取决于约束本身，与完整表的谱隙无关。

映射估计先固定模型族，因为有限个点对确定不了任意的光滑映射。刚体模型用质心消去平移，再由交叉矩阵的奇异值分解求旋转；允许反射的正交对齐用于比较重构坐标，带行列式修正的旋转对齐才对应刚体位姿。点重合或共线时交叉矩阵秩不足，最优旋转成族出现，零残差并不说明变换已被确定。仿射模型用增广坐标化为线性最小二乘，源点仿射张成定义域时参数唯一，所得的$\bm A$是否可逆则要另行核验。

射影模型在齐次坐标下只确定到一个非零倍数。四对处于一般位置的点使直接线性变换的核恰为一维，且对应可逆矩阵；点共线时齐次解可能全是奇异的，必须核验。含噪时先作坐标归一化，再用单位范数的齐次最小二乘求代数解；代数残差是以$c_i^2$加权的重投影误差，依赖于坐标表示，代数最优不等于几何最优。几种模型共同说明：残差小、方程有解与参数唯一、结果稳定是不同的结论，要分别核对。

本章的点之间要么有距离，要么有对应，对象始终是有限个点本身。第~\ref{chap:manifold-recovery}章把有限个点看成一个未知流形上的样本，研究怎样由它们恢复切方向、局部曲面、沿流形的距离以及同调，那里的距离图与样本复形都只用到本章意义下的成对距离。

